% Ensayo de Cornejo
El artículo de Anahiby Becerril, "The value of our personal data in the Big Data and the Internet of all Things Era", aborda un tema que, después de leerlo, uno no puede dejar de notar en la vida diaria: cada clic, cada ubicación y cada segundo que pasamos conectados a un dispositivo se convierte en un dato con valor monetario para las empresas de tecnología. Me parece un tema sumamente relevante porque, a diferencia de otros riesgos digitales más discutidos como el robo de contraseñas o el fraude, la explotación de nuestra información personal ocurre de manera silenciosa, dentro de términos y condiciones que casi nadie lee (incluyéndome a mí por supuesto). Mi postura es que la autora tiene razón al señalar la asimetría de poder entre empresas y usuarios, pero creo que su conclusión, que basta con "ser conscientes" del valor de nuestros datos para revertir el problema, se queda corta frente a la escala del fenómeno que ella misma describe.

Becerril construye su argumento apoyándose en cifras: mientras el Big Data y el IoT generan ingresos multimillonarios para compañías como Google y Facebook, estudios como los de Carrascal et al. o Staiano et al. (citados en la misma lectura) muestran que las personas le damos muy poco valor monetario a nuestra propia información. Esta separación entre lo que vale nuestra información para terceros y lo que nosotros creemos que vale es, para mí, el punto más interesante del texto. A partir de ahí, la autora plantea que el "precio" que pagamos por servicios aparentemente gratuitos es en realidad nuestra privacidad, y propone que herramientas como los Personal Data Stores o el proyecto The HAT, junto con regulaciones como el GDPR (de las pocas cosas que la Unión Europea sí está haciendo bien), son el camino hacia un modelo donde el individuo recupere el control de su información.

Aquí es donde me gusta traer a colación el trabajo de Jaron Lanier, uno de los autores que más disfruto sobre este tema. En su libro "Who Owns the Future?", Lanier no solo diagnostica el problema (él acuñó el término "siren servers" para describir a las plataformas que centralizan datos y poder), sino que propone algo más radical que la simple concientización: un sistema de micropagos donde, cada vez que nuestros datos generan valor económico, recibamos una compensación directa, en lugar de un servicio "gratuito" a cambio de renunciar a nuestra privacidad. Comparada con esta propuesta, la de Becerril se siente más tímida: pedir que los usuarios "sepan" cuánto valen sus datos no cambia el hecho de que, en la práctica, casi nadie tiene una opción real de negociar. Cuando instalamos una aplicación, la alternativa a aceptar los términos no es negociar mejores condiciones, sino simplemente no usar el servicio, lo cual la propia autora reconoce como una forma de exclusión social.

El objetivo central de Becerril, entiendo, no es convencernos de abandonar Internet, algo que ella misma reconoce como poco realista, sino persuadirnos de que la ignorancia sobre el valor de nuestra información es lo que permite que esta asimetría de poder se mantenga. La temática central del artículo es, entonces, la mercantilización de los datos personales, mientras que de forma lateral toca el crecimiento del IoT, la gobernanza de la información y el papel de la regulación. Esto conecta directamente con Fundamentos de Bases de Datos porque toda esta economía de datos depende, en última instancia, de cómo se diseñan y administran los sistemas que almacenan y procesan la información: la capacidad de las empresas para monetizar nuestros datos existe porque hay arquitecturas capaces de capturar, integrar y explotar volúmenes masivos de información en tiempo real. Como futuro profesional del área, esto me hace pensar que decisiones aparentemente técnicas, como qué campos se almacenan, cuánto tiempo se conservan o quién tiene permisos de acceso, tienen implicaciones éticas que van mucho más allá del diseño del esquema.

En conclusión, sostengo mi postura inicial: Becerril acierta en mostrar que existe una asimetría real entre lo que las empresas ganan con nuestros datos y lo que nosotros percibimos que valen, pero su propuesta de solución, basada en la concientización individual, me parece insuficiente sin mecanismos estructurales más fuertes, ya sea en forma de regulación o de modelos económicos alternativos como el que propone Lanier. Lo más valioso que me deja esta lectura es entender que el valor de mis datos no es abstracto sino cuantificable, y que como futuro diseñador de bases de datos tengo cierta responsabilidad sobre cómo esa información se estructura y protege. El texto también deja preguntas abiertas: ¿es posible construir un mercado de datos verdaderamente justo sin que termine replicando las mismas asimetrías que buscamos eliminar? Y si algún día pudiéramos vender directamente nuestra información, ¿tendríamos el poder de negociación necesario, o seguiríamos, otra vez, aceptando lo que las plataformas decidan ofrecernos?
